\subsection{Alcanos y alquenos}
Los \textbf{alcanos} son hidrocarburos saturados con enlaces sencillos carbono--carbono; los \textbf{alquenos} poseen al menos un doble enlace C=C. En los alquenos, el componente $\pi$ del doble enlace resulta pertinente para interpretar diversas adiciones. La estructura y las condiciones de reacción determinan el comportamiento de cada sistema: una categoría asignada a partir del nombre de una reacción no sustituye la comprobación del mecanismo molecular.

\subsection{Relación entre estructura y propiedades}
La masa molecular y la geometría de una molécula influyen en sus propiedades físicas. Para hidrocarburos apolares, la posibilidad de contacto intermolecular afecta las interacciones de dispersión. Así, dos isómeros con idéntica fórmula molecular pueden presentar puntos de ebullición distintos: la ramificación del esqueleto carbonado modifica la superficie de contacto efectiva. En esta investigación se examina particularmente la comparación entre butano e isobutano, sin extrapolar una única pareja de sustancias a una ley estadística.

\subsection{Interpretación de registros de reactividad}
Se separan dos roles documentales: en \emph{reacción} el hidrocarburo estudiado se registra como sustancia de partida y se caracteriza un producto principal; en \emph{obtención} el hidrocarburo es el compuesto de interés producido y se caracteriza un reactante evaluado. En ambos casos, las condiciones de temperatura, medio, solvente y catalizador son relevantes para describir el registro. El identificador \emph{CAS Reaction Number} localiza la reacción en la fuente documental, pero no reemplaza la identificación del compuesto o de su rol dentro de cada fila.
